\documentclass[11pt]{article}
\usepackage[a4paper,margin=2.6cm]{geometry}
\usepackage{amsmath,amssymb,amsthm}
\usepackage{booktabs}
\usepackage{array}
\usepackage{microtype}
\usepackage[hidelinks]{hyperref}
\usepackage{enumitem}
\usepackage{xcolor}
\usepackage{mdframed}

\newtheorem{theorem}{Theorem}
\theoremstyle{remark}
\newtheorem*{corollary}{Corollary}

\newmdenv[linewidth=0.8pt,linecolor=black!30,backgroundcolor=black!3,
          innertopmargin=8pt,innerbottommargin=8pt,skipabove=10pt,skipbelow=10pt]{caveat}

\newcommand{\FL}{\mathsf{FL}}
\newcommand{\RS}{\mathsf{RS}}
\newcommand{\RR}{\mathsf{RR}}
\newcommand{\Dd}{\mathsf{D}}
\newcommand{\sg}{\mathsf{sig}}
\newcommand{\afterr}{\mathsf{after}}
\newcommand{\before}{\mathsf{before}}

\title{\textbf{Negotiated exit from a deployed eUTxO escrow}\\[4pt]
\large Reachability, fee incidence, and the bargaining set of the Masumi \texttt{vested\_pay} validator}
\author{Brume}
\date{TOKEN2049 Origins, Cardano track \textbullet{} 6--7 October 2026}

\begin{document}
\maketitle

\begin{abstract}
\noindent A deployed Plutus V3 validator holds payments between autonomous agents on Cardano. We ask a question its
interfaces do not answer: from a given escrow, what can each party do, and which outcomes are reachable without a
third party.

We derive the validator's guard table from its source, then check it against the deployed bytecode by submitting
transactions. The main structural result is that, once the seller's deadline for delivering has passed, a seller's
concession is irreversible and removes arbitration for good. It leaves the buyer as the only party able to move the
escrow at all, and the only branch that takes value out of the script constrains no output.

Two consequences follow. Almost no split is self-enforcing, so almost none is reachable by incentives alone. And
every split becomes reachable by construction, through a transaction pair whose second leg is signed before the
first is submitted. We then price the bargaining set from measured arbitration history rather than from assumption.

We also state what the construction does not do. It removes a refusal. It does not remove a race.

All measurements are on Cardano mainnet and preprod, 6 October 2026. The two networks carry the identical script
hash, so every preprod execution interrogates the same bytes that hold the mainnet value.
\end{abstract}

\clearpage

\section{Setting and notation}

An escrow is a UTxO $u = (V, d)$ at the validator's script address. $V$ is a value: a family $(q_j)_{j \in J}$ of
quantities indexed by asset class $J$, lovelace included, which we write $j = L$. $d$ is an inline datum of sixteen
fields. We need nine of them.

\begin{center}
\small
\begin{tabular}{@{}lp{5.0cm}p{6.4cm}@{}}
\toprule
symbol & field & meaning \\
\midrule
$b, s$ & \texttt{buyer}, \texttt{seller} & payment credentials \\
$\sigma$ & \texttt{state} & $\sigma \in \Sigma = \{\FL, \RS, \RR, \Dd\}$ \\
$h$ & \texttt{result\_hash} & a byte string. $h = \varepsilon$ denotes empty \\
$c$ & \texttt{collateral\_return\_lovelace} & a lovelace floor returned to the buyer on one branch \\
$t_0$ & \texttt{unlock\_time} & the seller's lower gate on withdrawal \\
$t_1$ & \texttt{submit\_result\_time} & the seller's delivery deadline \\
$t_2$ & \texttt{external\_dispute\_unlock\_time} & the moment arbitration opens \\
$\theta_s, \theta_b$ & \texttt{seller\_cooldown\_time}, \texttt{buyer\_cooldown\_time} & per-party lower gates \\
$A, k$ & \texttt{admin\_vks}, \texttt{required\_admins\_multi\_sig} & the arbitration quorum, $k$ of $|A|$ \\
\bottomrule
\end{tabular}
\end{center}

\noindent A spend presents a redeemer $a \in \mathcal{A}$, where
\[
\mathcal{A} = \{a_0\,\mathsf{Withdraw},\; a_1\,\mathsf{SetRefundRequested},\; a_2\,\mathsf{UnSetRefundRequested},\;
a_3\,\mathsf{WithdrawRefund},
\]
\[
a_4\,\mathsf{WithdrawDisputed},\; a_5\,\mathsf{SubmitResult},\; a_6\,\mathsf{AuthorizeRefund}\}.
\]

The validator is then a partial labelled transition system. Each $a$ carries three independent guard families: a
signature predicate, a state predicate on $\sigma$ and $h$, and a temporal predicate on the transaction's validity
interval. We write $\afterr(t)$ for ``the interval starts after $t$'' and $\before(t)$ for ``the interval ends before
$t$''. Separately, each $a$ may impose output constraints, which are rules about where the spent value must go.

That separation is where the result of this paper lives. A redeemer can be tightly guarded about \emph{who} may call
it and say nothing at all about \emph{where the money goes}.

The protocol fee is $\varphi$ per mille. On mainnet $\varphi = 50$. We measured it at exactly $50.000$ per mille of
the token quantity in 8 of 8 observed $a_0$ spends, with the collateral return carrying zero tokens in all 8.

\section{The guard table}

Derived from the validator source, then checked by submitting transactions against the deployed bytecode. The four
redeemers that bear on the result are below; $a_1$, $a_2$ and $a_5$ move between states without exiting the contract
and are given where they are used, in the proof of Theorem~1.

\begin{center}
\small
\begin{tabular}{@{}llllp{1.7cm}p{5.0cm}@{}}
\toprule
$a$ & signature & states & on $h$ & time & output constraints \\
\midrule
$a_0$ & $\sg(s)$ & $\RS$ & $h \neq \varepsilon$ & $\afterr(t_0)$ &
  \textbf{yes.} Per asset $j$, at least $\lfloor q_j\varphi/1000 \rfloor$ to the fee address, and at least $c$
  lovelace to $b$, at datum-tagged outputs. The remainder is unconstrained \\
$a_3$ & $\sg(b)$ & $\FL, \RR$ & $h = \varepsilon$ & $\afterr(t_1)$ & \textbf{none} \\
$a_4$ & $k$ of $A$ & $\Dd$ & $h \neq \varepsilon$ & $\afterr(t_2)$ & \textbf{none} \\
$a_6$ & $\sg(s)$ & $\Dd$ & --- & $\afterr(\theta_s)$ &
  continuing output with $\sigma' = \RR$, $h' = \varepsilon$, value preserved \\
\bottomrule
\end{tabular}
\end{center}

\paragraph{Two remarks, and both matter for what follows.}
First, the arbitration branch $a_4$ constrains no output. The quorum may distribute the value as it chooses. In the
successor version of this validator the split is lifted into the redeemer itself, as two explicit amount fields. The
contract is not missing the notion of a negotiated split, and we do not claim it is.

Second, $a_3$ also constrains no output. In the successor version the source documents this as deliberate: a
payout-binding address may be set when the escrow is created, and when it is left unset the buyer signs the
destination. We are not the first to observe the property. We are the first to use it for a settlement between the
two parties.

\section{Arbitration closure}

\begin{theorem}
Let $u$ be an escrow with $\sigma = \Dd$ and $h \neq \varepsilon$, and suppose the present moment is past $t_1$, so
that no transaction can carry a validity interval ending before $t_1$. Let $u'$ be the continuing escrow produced by
a seller-signed $a_6$. Then the redeemers admissible on $u'$ are exactly $a_2$ and $a_3$. Both require $\sg(b)$, and
$a_3$ is the only one that removes value from the script. In particular $a_4$ is inadmissible on $u'$, and no
sequence of transitions from $u'$ restores it.
\end{theorem}

\paragraph{In plain terms.} Once the seller's delivery deadline has passed and the seller agrees to refund, the
arbitrators are locked out for good. From that moment only the buyer can move the escrow at all, and the branch that
takes the money out says nothing about where it goes.

\begin{proof}
By the $a_6$ output rule, $u'$ has $\sigma' = \RR$ and $h' = \varepsilon$. Take each redeemer in turn.

\begin{center}
\small
\begin{tabular}{@{}lp{11.4cm}@{}}
\toprule
$a$ & status on $u'$ \\
\midrule
$a_0$ & inadmissible: requires $\sigma = \RS$ \\
$a_1$ & inadmissible: its admissible states are $\FL$, $\RS$ and $\Dd$, and $\sigma' = \RR$ \\
$a_2$ & \textbf{admissible}, buyer-signed. It preserves $h$ and yields $\sigma = \FL$, because $h' = \varepsilon$ \\
$a_3$ & \textbf{admissible}, buyer-signed: $\sigma' \in \{\FL,\RR\}$, $h' = \varepsilon$, past $t_1$. It is the only
        branch here with no continuing output, so the only one that removes value \\
$a_4$ & inadmissible: requires $\sigma = \Dd$ \textbf{and} $h \neq \varepsilon$, and both are false \\
$a_5$ & inadmissible: its state guard is vacuous, so only time and $h$ can block it. It requires
        $\before(t_1)$, which no transaction can satisfy past $t_1$, \textbf{or} $h \neq \varepsilon$ together with
        $\before(t_2)$, and $h' = \varepsilon$ falsifies that \\
$a_6$ & inadmissible: requires $\sigma = \Dd$ \\
\bottomrule
\end{tabular}
\end{center}

\noindent Both admissible transitions are buyer-signed, so no party other than the buyer can move $u'$ at all.

It remains to show that no \emph{sequence} restores $a_4$. The only redeemer that writes a non-empty $h$ is $a_5$,
and the argument above blocks $a_5$ on \emph{every} state once $t_1$ is past, not only on $\RR$: its state guard is
vacuous, so the two temporal disjuncts are the whole of its guard, and the second needs the $h$ it would be writing
to be non-empty already. Hence $h = \varepsilon$ is invariant along every transition sequence from $u'$, including
through $\FL$ after $a_2$ and back through $a_1$. Since $a_4$ requires $h \neq \varepsilon$, it is unreachable.
\end{proof}

\paragraph{Why the hypothesis is not a technicality.} Before $t_1$ the theorem is false, and interestingly so. From
$\RR$ with $h = \varepsilon$, a seller-signed $a_5$ with $\before(t_1)$ writes a non-empty $h$ and, because the
current state is neither $\FL$ nor $\RS$, sets the new state to $\Dd$. That is exactly the pair $a_4$ needs, restored
in one step. Arbitration closure is a statement about escrows whose delivery deadline has passed. Every escrow we
measured on mainnet, and every one we settled on preprod, is past $t_1$; indeed $a_3$ itself cannot be called before
it.

\begin{corollary}[at source-reading strength only]
The same holds in version~2 of the validator. There $a_6$ writes a distinct state and still empties $h$, while $a_4$
still requires $\sigma = \Dd$ and $h \neq \varepsilon$. The result is not an artefact of one deployment. In version~2
the exit $a_3$ is unconstrained if and only if the optional payout-binding address is unset. Version~2 has never been
used on mainnet and we executed nothing against it, so every V2 statement here is \emph{the validator's source says}.
For V1 the refusals were observed on chain.
\end{corollary}

\paragraph{What Theorem 1 does not say.} It does not say the seller is cheated. A concession is voluntary, and the
four instances we observed on mainnet conceded a full refund. What it says is narrower and harder: after a concession
the value moves only if the buyer acts, and no third party can compel the outcome. Those four have been unmoved for
81 days.

\section{Fee incidence depends on the exit path}

Two orderings reach a settlement from $\Dd$.

\textbf{Path A.} The buyer plays $a_2$, so $\sigma \to \RS$. The seller then exits with $a_0$, paying the buyer a
share.

\textbf{Path B.} The seller plays $a_6$, so $\sigma \to \RR$ and $h \to \varepsilon$. The buyer then exits with
$a_3$, paying the seller a share.

Exiting through $a_0$ costs, per asset,
\[
\text{fee}^{A}_j = \left\lfloor \frac{q_j \varphi}{1000} \right\rfloor \le \frac{q_j\varphi}{1000},
\qquad\text{while}\qquad \text{fee}^{B}_j = 0 .
\]
At $\varphi = 50$ the pair keeps up to $5\,\%$ of each asset by settling through path B instead of path A, and
exactly $5\,\%$ whenever $1000$ divides $q_j\varphi$. This is a first-order term and it favours B without ambiguity.

The protective floor behaves the same way. Path A guarantees the buyer $c$ lovelace and no token. Path B guarantees
nothing at all. On the 61 mainnet escrows currently in $\Dd$, the contested quantity is overwhelmingly the token:
498.15 USDM against 54.31 ADA once collateral is removed. The one protective rule in the contract therefore protects
almost none of what is actually in dispute.

\section{Almost no split is self-enforcing}

An agreement assigns the first mover a share $\alpha_j \in [0,1]$ of each asset $j$, so the second mover is to
receive $(1-\alpha_j)q_j$. The second mover signs leg~2 and therefore holds a unilateral option: nothing in the
validator conditions admissibility on the agreement, because the agreement has no on-chain representation. The
question is when honouring it costs the second mover nothing.

\begin{theorem}
Fix a path and consider the second mover's choice at leg~2.
\begin{enumerate}[nosep,leftmargin=*]
\item \textbf{Path B.} The exit $a_3$ has no output constraints, so the buyer's best deviation pays them the whole
pot. Honouring is weakly optimal only if $\alpha_j = 0$ for every $j$.
\item \textbf{Path A.} The exit $a_0$ already compels the seller to send at least $c$ lovelace to the buyer and
$\text{fee}^{A}_j$ of each asset to the fee address. Honouring therefore costs the seller nothing exactly when the
agreement asks for no more than those mandatory outputs already deliver to the buyer, that is when
\[
\alpha_L q_L \le c \qquad\text{and}\qquad \alpha_j = 0 \ \text{ for every } j \neq L .
\]
\end{enumerate}
In particular no agreement that pays the first mover any token, or more lovelace than the collateral floor, is
self-enforcing on either path.
\end{theorem}

\paragraph{In plain terms.} Whoever signs second can always take everything the contract does not force them to give
up. On path B the contract forces nothing, so only a split that gives the first mover nothing survives. On path A it
forces a small lovelace floor, so the surviving splits are the ones already inside that floor.

\begin{proof}
Both cases compare two admissible transactions available to the second mover at leg~2 with identical inputs.

\emph{Path B.} Deviating, the buyer pays $q_j$ to themselves for every $j$; honouring, they pay
$(1-\alpha_j)q_j$ to themselves and $\alpha_j q_j$ to the seller. Both transactions satisfy $a_3$, which imposes no
output rule, so the deviation is admissible and weakly dominates unless $\alpha_j q_j = 0$ for every $j$.

\emph{Path A.} Both transactions must carry the tagged fee outputs and the $c$ lovelace to the buyer, so those
amounts are common to the two branches and cancel. What remains is the \emph{additional} transfer to the buyer the
agreement asks for, namely $\max(0, \alpha_L q_L - c)$ in lovelace and $\alpha_j q_j$ in each token. The deviation
withholds exactly that, and it is admissible because $a_0$ leaves the remainder unconstrained. So honouring is weakly
optimal precisely when that additional transfer is zero, which is the stated condition.
\end{proof}

The contract therefore makes every negotiated split worth negotiating unreachable by incentives alone. That is the
sharpest statement of the problem. It is also the statement of what has to be built.

\section{Reachability by construction}

Let $T_1$ be leg~1 and $T_2$ be leg~2. On Cardano a transaction identifier is a hash of its body. So
$\mathrm{id}(T_1)$ is fixed as soon as $T_1$ is built, before it is submitted. $T_2$ may therefore spend the output
$\mathrm{id}(T_1)\#0$, which does not yet exist, and may be signed against it.

The construction is this. Build $T_1$ and $T_2$ together. Have the second mover sign $T_2$. Only then does the first
mover sign and submit $T_1$, with $T_2$ following immediately.

\paragraph{Refusal is removed.} A signature on an unsubmitted transaction transfers nothing, so the second mover
gives up nothing by signing first. Once $T_2$ is signed there is no act of refusal left to perform.

\paragraph{Input ownership is not a detail.} The fee and collateral inputs of $T_2$ must belong to the conceding
party. If they belonged to the second mover, that party could invalidate $T_2$ at no cost by spending them elsewhere,
while keeping the option to take everything later.

\begin{caveat}
\textbf{What this does not provide.} Pre-signing removes a refusal. It does not remove a front-run. Once $T_1$ is
visible, the second mover may construct and broadcast a competing exit.

We submit both legs chained to one node. They landed in the same block in 5 runs of 5: blocks 5259678, 5259750,
5259751, 5259754 and 5259799, with leg~2 acknowledged 0.5 to 0.7\,s after leg~1 was sent. An adversary ran in four of
those five. In none of them did it obtain a clean window, so we report no front-run measurement and the solver prices
the event at its worst case.

The exposure is bounded below by the submitter's own gap between legs, measured at 249 to 296\,ms in the contended
runs, against an adversary whose first sighting came as early as 314\,ms. Those are the same order of magnitude. We
make no safety claim.
\end{caveat}

\section{The bargaining set, priced from observation}

Each party's reservation value is the value of not settling, which is the value of arbitration.

\paragraph{The seller's reservation value is zero, and that is measured.} Across all 120 lifetime $a_4$ transactions
on mainnet, the seller received nothing in any asset: 0 of 120. So $r_s = 0$, and any strictly positive share is
individually rational for the seller.

\paragraph{The buyer's side cannot be given a value, only a break-even.} The measured arbitration split is
$\pi_{\mathrm{ADA}} = 1.000$ and $\pi_{\mathrm{tok}} = 0.736$ of the pot: the fraction the buyer received, per asset,
in those 120 settlements. Model the arrival of the \emph{next} arbitration as a Poisson process of rate $\lambda$. No
$a_4$ transaction has occurred in $T = 313$ days, which gives the one-sided 95\,\% upper confidence bound
\[
e^{-\lambda T} = 0.05 \;\Longrightarrow\; \bar\lambda = \frac{\ln 20}{T} \approx \frac{1}{104\ \text{days}},
\]
the exponential form of the rule of three, since $\ln 20 = 2.996$. This bounds the rate from above and does nothing
else. The point estimate is zero and no lower bound exists. We fit only the silent window, not the lifetime history:
the 120 events occurred under a regime that visibly ended, and pooling them would assume the homogeneity the silence
is evidence against.

It is therefore not admissible to turn $\bar\lambda$ into a value of waiting and subtract it from the pot. There are
two reasons, and the second is the binding one.

\begin{enumerate}[leftmargin=*]
\item A buyer facing a \textbf{deadline} $H$ has a well-defined claim $\pi_j(1 - e^{-\lambda H})$. At $\bar\lambda$
that evaluates to 6.5\,\%, 25.0\,\% and 57.7\,\% of their arbitration payoff for $H = 7, 30, 90$ days. But it is an
upper bound computed from an upper bound, about a deadline we never observed.
\item A buyer who can \textbf{wait indefinitely}, against an offer that stays open, gives up nothing by waiting. Such
a buyer accepts less than the arbitration payoff only through a time preference, and nothing in the chain data
measures a discount rate. At a plausible 10\,\% per year, the same 95\,\% bound makes waiting worth about 97\,\% of
the ADA.
\end{enumerate}

So we publish break-evens, not values. Let $\beta_j$ be the share of asset $j$ offered \emph{to the buyer}, and
suppose the buyer has a deadline $H$. If $\beta_j \ge \pi_j$ the offer already matches the best arbitration has ever
paid, and accepting weakly dominates waiting at every rate. Otherwise $0 \le \beta_j < \pi_j$ and waiting beats
accepting if and only if
\[
\lambda \;>\; \lambda^{\star}(\beta_j, H) \;=\; -\frac{1}{H}\,\ln\!\left(1 - \frac{\beta_j}{\pi_j}\right) \;>\; 0 ,
\]
with the break-even horizon $H^{\star}$ following by inversion. The solver reports $\lambda^{\star}$ and $H^{\star}$
next to the measured bound $\bar\lambda$. The reader supplies the belief; we do not supply it for them.

The acceptable band is conditional and must always be quoted with its horizon. Writing it as the \emph{seller's}
share $1 - \beta_j$, a buyer with deadline $H$ accepts exactly
\[
\mathcal{B}_j(H) \;=\; \left[\, 0,\; 1 - \pi_j\left(1 - e^{-\bar\lambda H}\right) \right] ,
\]
non-empty for every $H$ because $\pi_j \le 1$. On path A the seller must additionally fund the mandatory outputs, so
the ceiling falls by $\varphi/1000$ in every asset and by a further $c/q_L$ in lovelace. \textbf{No split is called
rational in this paper without a stated horizon or a stated discount rate.}

One further term keeps settlement attractive to both parties even if arbitration were certain to arrive. In the 120
observed arbitrations, 26.4\,\% of the token pot reached neither party. We keep it as a solver term and do not
publish it as an output. The gross flows to admin-controlled addresses include those wallets' own change, so the only
defensible quantity is the pot minus what the parties received.

\section{What was executed}

Every structural claim above was checked against the deployed bytecode by submitting transactions on preprod. Its
script hash is byte-identical to mainnet's.

\paragraph{Six settlements, twelve legs, no admin key in any of them.} Each leg's CBOR was read back from the chain
and its hash checked. Across all twelve, every vkey witness hashes to either the buyer or the seller in that escrow's
own datum. None of the three deployed admin key hashes appears anywhere.

The clearest single example settles a token pot in block 5259570 through the file-drop path:

\begin{center}
\small
\begin{tabular}{@{}l@{}}
\toprule
\textbf{Leg 1}\quad $a_6$ \texttt{AuthorizeRefund} \\
{\scriptsize\ttfamily be1a2161b8c451309549265337893cc05bd4f75e9a7e3b971ccca86d7bcb4df8} \\
\addlinespace
\textbf{Leg 2}\quad $a_3$ \texttt{WithdrawRefund} \\
{\scriptsize\ttfamily ccb04dd233010d1e71ca0ebacb68cc83c28247e78e16b5f84a096389b2eab637} \\
\bottomrule
\end{tabular}
\end{center}

Leg~1 spends the escrow in $\Dd$ and returns the whole pot to the script in $\RR$, signed by the seller alone. Leg~2
spends \texttt{leg1\#0} and the value leaves the script: the buyer is the required signer, the seller co-signs
because their own UTxO pays the fee. The split is 12 tADA and 6 tUSDM to the buyer, 8 and 4 to the seller. Other
settlements in the same set run 70/30 and 25/75, so the construction is not tied to one ratio.

We also executed the admin pair on our \textbf{own} deployment of the same compiled code. There $a_4$ is accepted
before a concession and refused after it, with the state and hash guards named in the failure. And we executed
negative controls, in which the engine predicts a refusal, the transaction is submitted anyway, and the validator
refuses it.

Claims that rest on reading rather than execution are labelled as such throughout the repository. On the shared
mainnet script we are not an admin and never will be, so the admin branch there remains \emph{the validator's source
says}, never \emph{we observed}.

\begin{caveat}
\textbf{No transaction was ever built, evaluated or submitted against a mainnet escrow.} Mainnet is read-only in this
work.
\end{caveat}

\section{Limitations}

\begin{enumerate}[leftmargin=*]
\item \textbf{The front-run is unmeasured}, not measured and negative. Four contended attempts failed for transport
reasons, which we diagnose rather than hide.
\item \textbf{The arrival-rate bound is a bound.} $\bar\lambda$ is a 95\,\% upper limit from zero observations in the
silent window. It is not an estimate, and the true rate may be far lower.
\item \textbf{The dispute rate is not computable from our data.} The contract has 51{,}329 lifetime transactions, but
one job spends several, so the job denominator is unknown. We report counts, 120 arbitrated and 61 open, and no rate.
\item \textbf{The model assumes the parties already agree on the facts.} Nothing here decides who was right, and
nothing here could. The deliverable is private and only its hash is public. That is the premise, not a gap.
\item \textbf{The deployed V1 bytecode is not reproducible} from the current public source. Compiling with the
toolchain named in the committed blueprint reproduces the unapplied hash exactly, but the applied script is 6{,}003
bytes against 6{,}323 on chain. The V2 deployment reproduces byte-exactly, which establishes the method and isolates
the gap to a code difference rather than to our parameters.
\end{enumerate}

\end{document}
